\documentclass[11pt]{article}
\usepackage[margin=1in]{geometry}
\usepackage{amsmath,amssymb,amsthm}
\usepackage[hidelinks]{hyperref}
\setlength{\emergencystretch}{2em}
\setcounter{MaxMatrixCols}{15}
\newtheorem{theorem}{Theorem}
\newtheorem{proposition}[theorem]{Proposition}
\newtheorem{corollary}[theorem]{Corollary}
\newcommand{\F}{\mathbb F}
\newcommand{\mr}{\operatorname{mr}}
\title{All-characteristic reflection lifts:\\sixteen sharp five-bridge certificates}
\author{Internal research working proof}
\date{16 November 2026}
\begin{document}
\maketitle

The binary lift spaces in our retained sixteen-case panel give only
lower bound four. Changing coefficients to $\F_5$ gives lower bound five
for every case, matching five-seed bridge upper certificates. A separate
implementation verifies these equalities for the diagrams in the pinned
primary table. An independent named-census transfer now checks all sixteen
certificates directly, and a separate planar-graph constructor identifies
every case as a five-tangle Montesinos knot. Their MRC equalities were
already covered by Boileau--Zieschang (1985). These are sharp controls for
the method, not newly resolved cases. The full conjecture remains open.

This is an extension of the lift-counting method. It does not assert a
binary affine-symplectic realization of the new odd-characteristic lifts.
Uniform arguments below are written proofs, without formal verification
or external specialist review.

\section{The intrinsic bound over an arbitrary field}

Fix a knot group quotient $\phi:G_K\twoheadrightarrow S_r$, $r\ge3$,
taking meridians to transpositions. For any field $F$, put
\[
 E_r(F)=\{x\in F^r:\textstyle\sum_i x_i=0\},\qquad
 H_r(F)=E_r(F)\rtimes S_r.
\]
For $p<q$, write $u_{pq}=e_p-e_q$ and
\[
 s_{pq}^a=(a u_{pq},(pq)),\qquad a\in F.
\]
These elements are involutions in every characteristic: $(pq)u_{pq}=-u_{pq}$.
The collection $\mathcal C$ of these elements is conjugation invariant
in $H_r(F)$. Indeed, conjugating by $(b,\pi)$ changes the translation to
\[
 b+\pi(a u_{pq})-(\pi(pq)\pi^{-1})b,
\]
which is a multiple of the root of the projected transposition.

For meridional generators $m_1,\ldots,m_s$, write
$\phi(m_i)=(p_iq_i)$ with $p_i<q_i$. A lift with
$m_i\mapsto s_{p_iq_i}^{a_i}$ satisfies homogeneous linear relator
equations. Their solution space $Z_\phi(F)\subset F^s$ parametrizes all
such actual homomorphisms, without quotienting by conjugacy.

\begin{theorem}\label{thm:field}
The space $Z_\phi(F)$ is intrinsic to the specified representation
$\phi$, up to its coordinate description, and
\[
             \mr(K)\ \ge\ \dim_F Z_\phi(F).
\]
\end{theorem}
\begin{proof}
The split lift exists, and lift translations are cocycles, with their
linear structure given by addition of translation parts. Membership in
$\mathcal C$ on a meridian is preserved under conjugacy and inversion.
Thus the condition is intrinsic, independently of a meridional
presentation.

Let $g_1,\ldots,g_w$ be any meridional generating family. Above each
fixed transposition $\phi(g_j)$, an admissible image has exactly one
scalar parameter. Evaluating translations on these words therefore
gives a linear map $Z_\phi(F)\to F^w$. The map is injective: equality
on a generating family gives equality of homomorphisms. Consequently
$\dim_F Z_\phi(F)\le w$. For $F=\F_q$ the same proof is the counting
inequality $q^{\dim Z_\phi}\le q^w$.
\end{proof}

The same principle applies beyond symmetric groups. If
$\rho:G_K\to\operatorname{GL}(V)$ has
$\operatorname{rank}(\rho(\mu)-I)=1$, restrict affine lifts to those with
$z(\mu)\in\operatorname{im}(\rho(\mu)-I)$. This condition is invariant
under conjugation: a conjugating translation adds a vector in that same
image line. It is also preserved by inversion, since
$\operatorname{im}(A^{-1}-I)=\operatorname{im}(A-I)$. Such restricted
cocycles form a linear space, and evaluating on any $w$ meridional
generators injects it into $w$ one-dimensional fibres. Thus
$\mr(K)\ge\dim Z_{\rho}^{\mathrm{root}}$ over its coefficient field.
Unlike the involutory specialization, relator evaluation must retain
the distinction between positive and negative letters.

The argument belongs to the existing colouring and restricted-cocycle
framework. Clark--Elhamdadi--Saito--Yeatman define counts of lifts of a
fixed quandle colouring in Section 6 and apply total colouring counts to
bridge index in Section 7 \cite{cesy}. Here the conjugation quandles are
$\mathcal C\to\{\text{transpositions}\}$ and the fibre over a fixed base
colouring has $q^{\dim Z_\phi}$ elements. Evaluation on arbitrary
meridional generators supplies the meridional bound. No priority claim
is made for this counting principle. General cohomological meridional
bounds also appear in Gao--Gu--Li--Wang, Lemma 6.2 \cite{gglw}.

\section{One integer matrix covers all characteristics}

Compute every lifted relator in $\mathbb Z^r\rtimes S_r$, treating the
$a_i$ as independent variables. Its translation coordinates give an
integer matrix $M_\phi$ with $s$ columns. Since the generator images
are involutions, a negative letter has exactly the same lifted image
as its positive letter. In particular, unoriented over/under Gauss data
suffices for this calculation; no crossing-sign convention is guessed.

Suppose the nonzero Smith invariant factors are
$d_1\mid\cdots\mid d_t$, all positive. Unimodular integer row and column
operations remain invertible over every field. Hence
\[
 \dim Z_\phi(F)=
 \begin{cases}
 s-t,&\operatorname{char}F=0,\\
 s-\#\{i:p\nmid d_i\},&\operatorname{char}F=p.
 \end{cases}
\]

\begin{proposition}\label{prop:max}
If $u$ of the nonzero Smith factors equal one, then
\[
       \max_{F\text{ a field}}\dim_F Z_\phi(F)=s-u.
\]
If a nonunit factor exists, any prime divisor of the first nonunit
attains this maximum. Otherwise every characteristic attains it.
\end{proposition}
\begin{proof}
Unit factors never vanish in a field, giving the upper bound $s-u$.
A prime divisor of the first nonunit divides all subsequent factors by
the divisibility chain, so all and only those factors vanish. If every
nonzero factor is a unit, the rank is constant in every characteristic.
\end{proof}

Thus a failure for a fixed symmetric colouring can be excluded in all
characteristics exactly, without an arbitrarily chosen prime cutoff.
Our producer retains every elementary row or column swap, integer
multiple addition, and row negation. The independent verifier replays
these invertible steps and checks the resulting diagonal and its
divisibility chain; it does not run the producer's Smith algorithm.

\section{Sharp certificates on the retained five-seed panel}

The panel was frozen before the earlier binary search. It consists of
the first sixteen source-order rows with diagram Wirtinger number five
and no reported sharp finite Coxeter hit in the official BKM table
\cite{bkm}. The workbook is pinned at commit
{\small\texttt{74fe52e57de6f91988f63ac105cbc157534ad966}}; original bytes and
source-row numbers are retained. No cases were replaced after testing.

\begin{theorem}\label{thm:panel}
For each knot diagram labelled by the following names in that table,
meridional rank and bridge number both equal five:
\[
\begin{array}{rrrr}
15a1381&15a1383&15a1389&15a1416\\
15a1599&15a1617&15a1618&15a1619\\
15a7591&15a7715&15a7717&15a7733\\
15a7749&15a7752&15a7763&15a7816.
\end{array}
\]
The named-diagram transfer and the known Montesinos family identification
are checked below. Novelty of these knot cases is excluded by that audit.
\end{theorem}
\begin{proof}
For each case, a retained five-seed propagation trace colours all
fifteen Wirtinger arcs. The Wirtinger-number theorem \cite{wirt}
therefore gives $b(K)\le5$.

For each case, eight normalized surjective $S_4$ transposition
colourings have integer lift matrix with nonzero Smith factors
\[
                (\underbrace{1,\ldots,1}_{10},5).
\]
There are fifteen scalar variables. In characteristic five the matrix
has rank ten and its kernel has dimension five. Theorem
\ref{thm:field} gives $\mr(K)\ge5$. With $\mr(K)\le b(K)\le5$,
both equalities follow. Independently generated explicit kernel bases
and every upper-bound trace are retained in
\texttt{sharp\_five\_certificates.json}.
\end{proof}

For example, label the arcs of $15a1381$ by $0,\ldots,14$ in the
order fixed by its Gauss code. One certified projected colouring is
\[
\begin{split}
 &(01),(01),(01),(01),(01),(02),(03),(23),\\
 &(23),(02),(03),(23),(02),(12),(01).
\end{split}
\]
Its five seeds are $\{0,1,5,6,7\}$. A basis of the lift kernel over
$\F_5$, written in the same arc order, is
\[
\setlength{\arraycolsep}{3pt}
\begin{pmatrix}
2&3&3&1&0&0&0&0&0&0&0&0&0&0&0\\
0&0&0&0&0&0&0&4&4&1&0&0&0&0&0\\
0&0&0&0&0&0&1&1&1&0&1&1&0&0&0\\
0&0&0&0&0&1&1&1&1&0&1&0&1&1&0\\
4&3&3&0&1&1&1&1&1&0&1&0&1&0&1
\end{pmatrix}.
\]
The five columns at the seed positions are independent. These vectors
therefore already certify five independent actual lifts, without
needing the full all-characteristic search to prove this lower bound.

The complete normalized distributions are identical on all sixteen
cases. For $S_3$, six colourings have factors $1^{12},30$ and two have
$1^{11},10$; their maximum bounds are respectively three and four.
For $S_4$, twelve have $1^{11},15$, four have $1^{11},10$, and eight
have $1^{10},5$; their maximum bounds are four, four, and five.
In characteristic two all $S_4$ bounds are three or four, as in the
earlier receipt. The new sharpness specifically requires characteristic
five. Individual image surjectivity beyond the base $S_4$ is not needed.

\section{Named-diagram transfer and the prior-family audit}

SnapPy's independently packaged Hoste--Thistlethwaite census supplies
planar diagrams for all sixteen names. The initial over/under Gauss
walks agree with the workbook up to traversal choices and global mirror.
The standard-library verifier then transfers every explicit $\F_5$
kernel basis and seed trace to the census diagram itself and checks
every crossing there. Thus named-knot certification no longer depends
only on accepting the workbook's knot labels.

A frozen family audit first tested 112 five-tangle Montesinos diagrams
suggested by the crossing budget, without replacing any target. Seven
matched exactly. A separately frozen second stage retained all those
fraction tuples and tested all $2^5$ individual planar half-turn choices,
3584 diagrams in total. The rational-tangle flipping lemma
\cite[Lemma 2(iii)]{kl} preserves each fraction under these half-turns.
All sixteen then matched. No numerical isometry or invariant-only test
was accepted.

A separate implementation constructs the rational tangles from paired
endpoints using integer twists, inversion, side-by-side sum and numerator
closure. It checks each continued fraction exactly over $\mathbb Q$.
An explicit bijection of crossing half-edges checks the generated map
against the retained candidate and then the named census diagram.
It preserves edge pairing and cyclic rotation, allowing only a single
global reversal of planar orientation and a single global crossing
switch. These transformations correspond to ambient reflections; MRC
and both ranks are invariant under mirror. A control switching just one
crossing is rejected. The witnesses are retained in
\texttt{diagram\_audit\_verification.json}.

For example, up to mirror,
\[
 15a1381=M\left(0;\frac12,\frac35,\frac23,\frac23,\frac23\right).
\]
Every identified presentation has five reduced, noninteger fractions,
with denominator multiset $\{2,3,3,3,5\}$. All sixteen therefore fall
under the classical Montesinos theorem $\mr=b=r$ for $r\ge3$ branches
\cite{bz,musick}. The original theorem statement is restated explicitly
in Musick's Theorem 1; the entire 1985 proof has not been audited here.
The F5 certificates give another direct verification of these known
equalities, not a new knot-class theorem.

\section{An unlimited-radical obstruction and a current uncertain panel}

\begin{proposition}\label{prop:radical}
Suppose a knot has a transitive affine representation on a binary vector
space $V$, whose meridians have the common-form transvection shape
\[
 T_{u,a}(x)=x+(B(u,x)+a)u,
\]
where the alternating form $B$ has rank two. Then the knot group has a
surjective $S_3$ transposition representation. This holds for any
dimension of the radical of $B$.
\end{proposition}
\begin{proof}
Let $R=\operatorname{rad}B$. The linear parts induce a representation
on $W=V/R$, with image in $\operatorname{Sp}(W)=\operatorname{Sp}(2,2)
\cong S_3$. Transitivity forces the meridian directions to span $V$:
every generator changes a point by a vector in the span of these
directions, so its cosets would otherwise be invariant. Their projected
directions therefore span the two-dimensional $W$.

At least two distinct nonzero projected directions occur. Their linear
parts are distinct transvections in $\operatorname{Sp}(2,2)$, hence
distinct transpositions which generate $S_3$. Since meridians are
conjugate, none can have trivial projected linear part when another
has nontrivial part. The resulting quotient has the required meridian
images.
\end{proof}

We separately froze the first sixteen names in the 293-case
\texttt{still\_open.json} list of the October 2026 release
\cite{bhat}, pinned at commit
\path{06cbbb7e2737e1a11085b64e0ead021d63c550c3}. Joining names to
the BKM workbook gives independent four-seed diagram upper certificates.
Exhaustive normalized assignments give no surjective $S_3$ or $S_4$
transposition colouring for any of these sixteen diagrams. Independent
permutation enumeration reproduces this absence. These are
\[
\begin{array}{rrrr}
13a129&13a135&13a136&13a137\\
13a194&13a196&13a197&13a198\\
13a233&13a234&13a239&13a242\\
13a243&13a244&13a253&13a265.
\end{array}
\]
Consequently no $S_3$/$S_4$ lift strategy over any field applies to
them. The companion \path{BRANCHED_COVER_LIFT_BOUNDS.tex} completes
the exclusion: a four-meridian generating family permits a connected
transposition quotient only in degrees $r\le5$. An exhaustive $S_5$
search also gives no base, while the unique $S_2$ base has maximum
lift dimension two in every characteristic, by its integer Smith form.
Thus this symmetric root-lift method cannot give a sharp four-bound
on any of these targets, in any degree or field.

The original binary common-form affine criterion is also excluded in
\emph{every} dimension: its bound forces $k+1\le\mr\le b\le4$, so
$k\le3$ and every nonzero alternating form has rank two. Proposition
\ref{prop:radical} would then force the absent $S_3$ quotient. The
earlier suggestion that higher-rank binary forms remain possible on
this panel is withdrawn. Other quotient types and lower-bound methods
remain possible; these exclusions do not resolve the targets' MRC status.

We also froze these same sixteen targets with their oriented PD data
from KnotInfo 2026.11.2 and tested positive transvections in
$\operatorname{SL}_2(\F_p)$ for $p=3,5,7$. For
$u=(a,b)\ne0$, the base matrix is
\[
 A_u=\begin{pmatrix}1-ab&a^2\\-b^2&1+ab\end{pmatrix},
 \qquad u\sim-u.
\]
Restricted affine lifts take a meridian to $(c u,A_u)$, $c\in\F_p$,
so the preceding general root-space bound applies. Four-seed propagation
certifies the upper bounds. Complete normalized noncommuting assignments
in the stated conjugacy classes have only lift dimensions two or three;
none has dimension four. Parallel root families give commuting affine
images, which must all coincide since knot meridians are conjugate;
their lift space has dimension one and cannot be sharp here.
A separate replay enumerates the conjugacy class
from all determinant-one matrices, uses the seed trace for propagation,
and evaluates dense three-by-three affine matrices to reconstruct the
lift equations. It reproduces every distribution. The PD parser alone
is shared. These negative results exclude only the stated finite fields
and classes; they do not exclude other primes or a general parabolic
strategy. The retained receipt is
\texttt{parabolic\_panel\_verification.json}.

\section{Verification and remaining gates}

The new separate verifier checks all 512 lift matrices in the earlier
panel by reverse affine point evaluation, replays integer elementary
operations, and cross-checks dimensions by independent elimination in
characteristics $2,3,5,7,11$. It enumerates projected permutation groups
and every normalized seed assignment independently. Smith controls
cover all 625 two-by-two matrices with entries in $\{-2,-1,0,1,2\}$,
using gcds of minors as a separate oracle. Both sixteen-case upper
traces are checked. The earlier independent workbook-reader replay is
retained. These are finite checks supporting the stated certificates,
not a formal verification of the uniform proofs.

We also retained the October release's 243 matrix certificates and
replayed all of them with a separate standard-library checker: column
rank elimination checks invertibility, rank-one differences, and zero
common fixed space; direct multiplication checks the signed Wirtinger
relations without matrix inversion; overstrand identifications check
the four-seed upper certificates. All pass. Changed matrix entries,
three seeds, and identity matrices are rejected. All 243 source diagrams
also have exact crossing-preserving planar-map matches in the separately
packaged HT census. KnotInfo releases 2026.10.5 and 2026.11.2 each have
exactly the same 536 thirteen-crossing $[3,4]$ entries as the author's
243-plus-293 partition; all 243 published PD records agree literally
with both snapshots. Original archive and CSV hashes and the selected
source columns are retained. This checks those archived tables, not
every subsequent live database update. The full paper remains unaudited
and unrefereed; credit for these 243 cases remains with its author.

The named-diagram and prior-family gates for the sixteen five-bridge
cases are complete: all are known Montesinos cases. A no-hit in the
finite Coxeter table was not evidence that MRC was previously open.
The current thirteen-crossing panel requires a different lower-bound
mechanism or a quotient outside the excluded symmetric root-lift route.
The companion branched-cover proof identifies lift nullity as
$r-1+\dim H_1(Y_\phi;F)$ and supplies a nonabelian cover-rank gate.
Its separate $A_5$ core-colouring count gives no sharp certificate on
the same sixteen targets. The full conjecture and the outstanding
source and recent-proof audits remain open.

\begin{thebibliography}{99}
\bibitem{bz} Boileau--Zieschang, Nombre de ponts et generateurs meridiens
des entrelacs de Montesinos (1985), Theorem 1.1.
\url{https://doi.org/10.1007/BF02567413}.
\bibitem{musick} Musick, Minimal bridge projections for 11-crossing
prime knots (2012), Theorem 1. \url{https://arxiv.org/abs/1208.4233v3}.
\bibitem{kl} Kauffman--Lambropoulou, On the classification of rational
tangles (2004), Theorem 1 and Lemma 2.
\url{https://arxiv.org/abs/math/0311499v2}.
\bibitem{cesy} Clark--Elhamdadi--Saito--Yeatman, Quandle Colorings of
Knots and Applications, J. Knot Theory Ramifications 23 (2014), 1450035,
Sections 6 and 7. \url{https://arxiv.org/abs/1312.3307}.
\bibitem{gglw} Gao--Gu--Li--Wang, Meridional ranks of iterated Whitehead
doubles (2026), Lemma 6.2. \url{https://arxiv.org/abs/2611.15673v1}.
\bibitem{bkm} Blair--Kjuchukova--Morrison, Coxeter quotients of knot
groups through 16 crossings (2024).
\url{https://arxiv.org/abs/2208.09032}.
Official data: \url{https://github.com/ThisSentenceIsALie/Wirt_Hm}.
\bibitem{wirt} Blair--Kjuchukova--Velazquez--Villanueva, Wirtinger
systems of generators of knot groups (2020).
\url{https://archive.intlpress.com/site/pub/files/_fulltext/journals/cag/2020/0028/0002/CAG-2020-0028-0002-a002.pdf}.
\bibitem{bhat} Bhattacharjee, The bridge index and the nonorientable
four-genus of knots with at most thirteen crossings, unrefereed release
1.3.0 (2026). \url{https://doi.org/10.5281/zenodo.23561157}.
Official data: \url{https://github.com/anaxia-research/conjecture-to-choose}.
\end{thebibliography}
\end{document}
